\section{RL\_AGENT/ems\_offline\_RL\_agent.py}

\subsection*{Scopo}
Runner operativo per training/eval/test RL offline, con gestione output, callback, log e TensorBoard.

\subsection*{Responsabilita}
\begin{itemize}[leftmargin=1.2cm]
\item Parsing argomenti CLI e selezione modalita (\texttt{train/eval/test/rollout}).
\item Override algoritmo da CLI con \texttt{--algorithm} (\texttt{ppo} o \texttt{sac}),
  applicato al campo \texttt{rl.algorithm} prima della build/load del modello.
\item Selezione dataset per modalita e validazione start/steps.
\item Split train/eval su finestre temporali dedicate:
  supporta \texttt{rl.train\_split\_start\_step}/\texttt{rl.train\_split\_end\_step} e
  \texttt{rl.eval\_split\_start\_step}/\texttt{rl.eval\_split\_end\_step}; il runner
  applica i limiti prima della creazione dell'env.
\item Compatibilita dataset-mode: supporta sia dataset CSV sia
  \texttt{scenario.dataset\_mode: pymgrid\_bundle} senza forzare path CSV su \texttt{eval/test}
  quando i dati arrivano dallo scenario.
\item Setup callback e logging (\texttt{TensorboardLogger}, checkpoint, metriche episodio).
\item Diagnostica reward/stato coerente con naming \texttt{soe\_*} (non \texttt{soc\_*})
  per violazioni limite e stato energetico osservato.
\item Export transition history robusto: plot/salvataggio leggono prima dal BMS e
  usano il transition model solo come fallback legacy.
\item Export risultati e grafici finali.
\item Nel \texttt{final\_report.txt} stampa lo stato \texttt{Degradazione SoH : ON/OFF}
  usando la configurazione batteria del simulatore.
\end{itemize}

\subsection*{Dipendenze}
\begin{itemize}[leftmargin=1.2cm]
\item Usa \texttt{EMS\_RL\_agent.py} per env e agente.
\item Usa \texttt{tb\_logger.py} per logging avanzato.
\end{itemize}

\subsection*{Note manutenzione}
Quando si cambiano flag CLI, policy di selezione dataset o percorsi output,
aggiornare script sweep/test RL e documentazione operativa.
Le coppie \texttt{*\_split\_start\_step}/\texttt{*\_split\_end\_step} sono riferite
all'indice globale del dataset; gli \texttt{rl.train\_start\_step}/\texttt{rl.eval\_start\_step}
vengono ricondotti automaticamente all'indice locale della finestra selezionata.
